%! Systems Involving Quadratics — Unit 3, Lesson 3.6 — Lesson Plan
\documentclass[10pt]{article}
\usepackage{algebra2-article}
\usepackage{algebra2-boxes}
\usepackage{graphicx}   % -article does NOT load graphicx; needed for thumbnails/screenshots
\graphicspath{{images/}}
\newcommand{\TallMath}[1]{$\displaystyle #1\rule[-1.4em]{0pt}{3.2em}$}
\newcommand{\CourseName}{Algebra 2: Shepherd}
\newcommand{\SchoolYear}{2026--2027}
\newcommand{\MeetingLength}{55 minutes}
\newcommand{\UnitNumberName}{Unit 3: Quadratic Functions \quad}
\newcommand{\LessonNumberName}{Lesson 3.6: Systems Involving Quadratics}

\begin{document}
\begin{center}
    {\Large\bfseries \CourseName: \SchoolYear \\
    {\small\bfseries \UnitNumberName \LessonNumberName}}
\end{center}

\begin{tcolorbox}[colback=forestbg, colframe=forest, boxrule=0.9pt, arc=2mm,
    left=2mm, right=2mm, top=1.5mm, bottom=1.5mm]
    \textbf{Primary Objective:} Students extend equation-solving from a single quadratic (Lessons 3.2--3.5) to
    a \emph{system} of two equations in which at least one is quadratic. They learn that a \textbf{solution} is
    an ordered pair satisfying \emph{both} equations --- graphically, a \textbf{point of intersection} of the
    two graphs --- and that a \textbf{linear--quadratic} system (line \& parabola) or
    \textbf{quadratic--quadratic} system (two parabolas) can have \textbf{0, 1, or 2} solutions. They solve
    algebraically by \textbf{substitution} (and, for quadratic--quadratic, by \textbf{elimination}), the key
    move being that substitution \emph{collapses the system into a single quadratic equation} they already know
    how to solve --- by factoring or the quadratic formula (Lesson 3.5) --- after which they back-substitute to
    recover the paired coordinate. The \textbf{discriminant} of that collapsed equation predicts the number of
    intersection points (2 / 1 / 0), tying today's ``how many solutions'' directly back to Lesson 3.5. Students
    verify solutions algebraically and graphically and model a contextual situation (break-even, meeting paths).
    \\[2pt]
    \textbf{Standards (2023 VA SOL):} \textbf{A2.EI.3c} (solve a linear--quadratic and quadratic--quadratic
    system algebraically and graphically, including in context), \textbf{A2.EI.3b} (determine the number of
    solutions), \textbf{A2.EI.3a} (create such a system to model a contextual situation), and \textbf{A2.EI.3d}
    (verify solutions algebraically, graphically, and with technology; interpret in context). Builds directly on
    solving a single quadratic by factoring (\textbf{A2.EI.2b}, Lesson 3.2) and the quadratic formula /
    discriminant (\textbf{A2.EI.2b}, \textbf{A2.F.2d}, Lesson 3.5).
\end{tcolorbox}

\renewcommand{\arraystretch}{1.4}
\begin{skillbox}[Priority Ideas \& Skills]{goldbox}
\begin{tabularx}{\linewidth}{@{}X|X@{}}
\textbf{Priority Ideas \& Skills} & \textbf{Key Understandings (the ``why'')} \\[2pt]
\hline
\vspace{-6pt}
\begin{itemize}[leftmargin=1.1em, nosep, after=\vspace{2pt}]
  \item Read the solution(s) of a system as the intersection point(s) of two pre-drawn graphs.
  \item Solve a linear--quadratic system by \textbf{substitution}: substitute the line into the parabola.
  \item Solve a quadratic--quadratic system by setting the expressions equal (or eliminating).
  \item Recognize that either method collapses the system into \emph{one} quadratic equation to solve.
  \item Determine the number of solutions (2 / 1 / 0) --- and connect it to the collapsed equation's discriminant.
  \item Verify a solution in \emph{both} equations; model and interpret a contextual system.
\end{itemize}
&
\vspace{-6pt}
\begin{itemize}[leftmargin=1.1em, nosep, after=\vspace{2pt}]
  \item A solution must satisfy \emph{both} equations at once --- that is exactly what ``intersection'' means.
  \item Substitution replaces $y$ so the two-variable system becomes a familiar one-variable quadratic.
  \item ``Set them equal'' works because at an intersection the two $y$-values agree.
  \item Every solving path funnels back to Lessons 3.2--3.5 --- there is no genuinely new equation to solve.
  \item A line can miss, touch, or cross a parabola; the discriminant counts the crossings before you solve.
  \item An algebraic answer is only a real solution if it checks in \emph{both} equations and makes sense in context.
\end{itemize}
\end{tabularx}
\end{skillbox}

\begin{skillbox}[{Vocabulary, Concepts, \& Theorems}]{sky}
\renewcommand{\arraystretch}{1.3}
\begin{tabularx}{\linewidth}{@{}l X@{}}
\textbf{System of equations} & Two or more equations in the same variables, considered together. \\
\textbf{Solution of a system} & An ordered pair $(x,y)$ that makes \emph{every} equation true at once. \\
\textbf{Linear--quadratic system} & One line and one parabola; graphs meet in $0$, $1$, or $2$ points. \\
\textbf{Quadratic--quadratic system} & Two parabolas; meet in $0$, $1$, or $2$ points (or coincide). \\
\textbf{Substitution} & Replace $y$ (from the linear equation) into the quadratic to get one equation in $x$. \\
\textbf{Elimination} & Subtract two ``$y=\dots$'' equations so the $y$'s cancel, leaving one equation in $x$. \\
\textbf{Point of intersection} & Where two graphs cross --- the graphical picture of a solution. \\
\end{tabularx}
\end{skillbox}

\begin{fixedskillbox}[Activate Prior Knowledge \& Spiral Review]{forestbg}
    The Warm-Up rehearses the three moves substitution will demand: \textbf{(1)} \emph{solving a single quadratic}
    $x^2-3x=0$ and $x^2-2x-8=0$ by factoring (Lesson 3.2), because that is exactly what a system collapses to;
    \textbf{(2)} \emph{evaluating a discriminant} to count real roots (Lesson 3.5), previewing ``how many
    intersection points''; and \textbf{(3)} \emph{checking whether an ordered pair satisfies an equation}
    --- the habit that becomes ``verify in both equations.'' Debrief the check last: ``a point on the parabola
    isn't a solution of the \emph{system} unless it's also on the line.''
\end{fixedskillbox}

\begin{skillbox}[Hook]{forestbg}
    Project the unit anchor parabola $y=x^2-2x-3$ with three horizontal lines drawn across it --- $y=5$ (cuts it
    twice), $y=-4$ (just touches the vertex), and $y=-6$ (misses entirely below). Ask: ``How many points does
    each line share with the parabola?'' Students count $2$, $1$, $0$ straight off the graph. Then the pivot:
    ``Those counts are the number of \emph{solutions} of a \emph{system}. And notice --- $y=-4$ touching once is
    the same `one repeated root' story as a zero discriminant from Lesson 3.5.'' Today's job is to find those
    intersection points \emph{exactly}, with algebra, when the graph won't hand you clean coordinates.
\end{skillbox}

\begin{skillbox}[Lesson]{forestbg}
\begin{multicols}{2}
\textbf{1. What a solution is.} A \emph{system} pairs two equations; a \emph{solution} is an ordered pair making
\emph{both} true --- graphically, a \textbf{point of intersection}. Show the anchor parabola $y=x^2-2x-3$ with the
line $y=x-3$; the graphs cross at $(0,-3)$ and $(3,0)$. Verify $(3,0)$ in \emph{both} equations to model the check.

\textbf{2. Substitution (linear--quadratic).} Both equations give $y$, so set the parabola's $y$ equal to the
line's $y$: $x^2-2x-3=x-3$. The $y$ disappears and a \emph{single quadratic} remains: $x^2-3x=0\Rightarrow
x(x-3)=0\Rightarrow x=0,3$. Back-substitute into the \emph{line} (easier) to get $y$: $(0,-3),(3,0)$. Post the
protocol: isolate $y$ $\to$ substitute $\to$ solve the one quadratic $\to$ back-substitute $\to$ write pairs $\to$ check.

\columnbreak

\textbf{3. How many solutions?} Return to the hook's three lines against $y=x^2-2x-3$: $y=5$ gives
$x^2-2x-8=0$ ($D=36>0$, \emph{two} points); $y=-4$ gives $x^2-2x+1=0$ ($D=0$, \emph{one} --- tangent);
$y=-6$ gives $x^2-2x+3=0$ ($D=-8<0$, \emph{none}). \emph{The discriminant of the collapsed quadratic counts the
intersections} --- Lesson 3.5, reused.

\textbf{4. Quadratic--quadratic.} Two parabolas: set them equal. $y=x^2$ and $y=-x^2+8$ give $x^2=-x^2+8\Rightarrow
2x^2=8\Rightarrow x=\pm2$: points $(\pm2,4)$. Warn: if the $x^2$ terms are \emph{identical} they cancel, leaving a
\emph{linear} equation and at most one solution (e.g.\ $y=x^2+1$, $y=x^2-2x+5\Rightarrow x=2$).

\textbf{5. Model.} Break-even: revenue $y=-x^2+24x$ meets cost $y=8x+48$ at $x=4$ and $x=12$ (hundreds of units);
between them revenue beats cost. Create, solve, verify, interpret (A2.EI.3a/c/d).
\end{multicols}
\end{skillbox}

\begin{skillbox}[Explicit Instruction: Substitution for a Linear--Quadratic System]{forestbg}
\begin{tabularx}{\linewidth}{@{}p{0.46\linewidth} X@{}}
\textbf{Steps (post and refer back):}
\begin{enumerate}[leftmargin=1.4em, nosep, itemsep=2pt]
  \item Solve the \emph{linear} equation for $y$ (or $x$).
  \item Substitute into the \emph{quadratic} --- one equation in one variable.
  \item Write it in standard form and solve (factor or formula).
  \item Back-substitute each $x$ into the \emph{line} to get $y$.
  \item Write each solution as an ordered pair $(x,y)$.
  \item Check each pair in \emph{both} original equations.
\end{enumerate}
&
\textbf{Worked example.} Solve $\;y=x^2+1,\ \ y=2x+4$.
\begin{itemize}[leftmargin=1.2em, nosep, itemsep=1pt]
  \item Substitute: $x^2+1=2x+4$.
  \item Standard form: $x^2-2x-3=0$.
  \item Factor: $(x-3)(x+1)=0\Rightarrow x=3,-1$.
  \item Back-sub into $y=2x+4$: $x=3\!\to\!10$, $x=-1\!\to\!2$.
  \item Solutions: $(3,10)$ and $(-1,2)$.
  \item Check $(3,10)$: $3^2+1=10$ \checkmark, $2(3)+4=10$ \checkmark.
\end{itemize}
\end{tabularx}
\end{skillbox}

\begin{skillbox}[Active Monitoring]{redbox}
Circulate during the notes practice and Tier work. Watch for and correct:
\begin{itemize}[leftmargin=1.2em, nosep]
  \item finding $x$ but \emph{stopping} --- forgetting to back-substitute for $y$ (a solution is a \emph{pair});
  \item back-substituting into the harder equation, or into only one, so the two coordinates don't match;
  \item sign slips when moving every term to one side to reach standard form $ax^2+bx+c=0$;
  \item pairing the $x$'s and $y$'s wrongly (e.g.\ crossing $x=3$ with the other point's $y$);
  \item on quadratic--quadratic, dropping a solution when the $x^2$ terms \emph{don't} cancel (it's still quadratic);
  \item calling a point that lies on one graph a ``solution'' without checking it on the \emph{other}.
\end{itemize}
Cold-call prompts: ``Both equations equal $y$ --- so what can you set equal?'' ``You found $x=3$; are you done?''
``The collapsed equation has discriminant $0$ --- how many times do the graphs meet?'' ``Which original equation is
easier to back-substitute into?''
\end{skillbox}

\begin{skillbox}[Group Work \& Differentiation]{redbox}
\textbf{Where Do They Meet?} activity (tiered; mirrors the notes).
\begin{multicols}{3}
\textbf{Tier R --- Remediate.} Read solutions of pre-drawn linear--quadratic and quadratic--quadratic systems as
intersection points; count $0/1/2$ solutions; verify a given ordered pair in \emph{both} equations.

\columnbreak
\textbf{Tier A --- Approaching.} Solve linear--quadratic systems by substitution (factorable and one
formula/irrational case) and a quadratic--quadratic system by setting equal; write ordered-pair solutions.

\columnbreak
\textbf{Tier E --- Extension.} Model a break-even/meeting-paths context (create, solve, interpret); use a
discriminant to find the value of a parameter making a line \emph{tangent} to a parabola (exactly one solution).
\end{multicols}
\end{skillbox}

\begin{skillbox}[Individual Work \& Assessment]{redbox}
\textbf{Exit Ticket} (independent, no notes): solve one linear--quadratic system by substitution (integer
solutions, two points); state the number of solutions of a second system and justify it; and an
\textbf{SOL-style multiple-choice} item on the number of solutions of a linear--quadratic system. Collect and sort
into ``got it / found $x$ but no $y$ / mismatched coordinates / didn't check both.'' to decide whether Lesson 3.7
opens with a re-teach.
\end{skillbox}

\begin{skillbox}[Reinforcement \& Extension]{goldbox}
\textbf{Homework:} verify ordered pairs in both equations; solve linear--quadratic systems by substitution (real
rational and irrational); solve a quadratic--quadratic system by setting equal / eliminating; determine the number
of solutions from a discriminant and match to a graph; and a short ``explain the error'' item on forgetting to
back-substitute. \textbf{Extension:} a break-even model and a find-the-parameter (tangent line) discriminant
problem. \textbf{Preview:} Lesson~3.7, \emph{Modeling with Quadratics} --- projectile, area, and optimization
problems where the quadratic model itself (and its vertex) carries the meaning.
\end{skillbox}
\end{document}
